% =====================================================================
% introduction.tex  --  "Polarization Defusing" project
% Draft Introduction. Hands off into % =====================================================================
% background.tex  --  "Polarization Defusing"  (LAB-ONLY related work)
% At this stage the related work reviews ONLY papers from Ray LC's lab.
% Requires in preamble: \usepackage{booktabs} \usepackage{amssymb} \usepackage{tabularx}
% Bibliography keys live in refs.bib.
% =====================================================================

\newcommand{\yes}{\checkmark}
\newcommand{\no}{--}
\newcommand{\prt}{$\sim$}   % partial / indirect

\section{Background and Related Work}
\label{sec:background}

Our project asks how to \emph{defuse} polarization in large interpersonal
online groups---recruiting already-polarized communities and intervening in
their everyday chat medium (forum threads, group messaging, Discord-style
channels) with a population of generative-AI agents that exert \emph{social
influence} rather than top-down moderation. Rather than survey the field
broadly, we position this project against our group's own prior work, which has
independently developed each ingredient the project must combine: (i)~studies of
how people express opinion and form communities on everyday online platforms,
(ii)~systems in which LLM and multi-agent technologies steer opinion or
behavior, and (iii)~conversational and playful designs that shift attitudes.
Table~\ref{tab:rw} maps representative lab systems against the design
commitments of the present project, and the three subsections below unpack each
line in turn. Read together, they show that while our group has built every
piece, no prior project has assembled them into a multi-agent intervention for
defusing polarization---precisely the gap we target.

% ---------------------------------------------------------------------
\begin{table*}[t]
  \centering
  \footnotesize
  \setlength{\tabcolsep}{4pt}
  \renewcommand{\arraystretch}{1.2}
  \caption{Map of prior \emph{lab} work against the design commitments of
    \emph{Polarization Defusing}. \yes~= central, \prt~= partial/indirect,
    \no~= absent. ``Target'' marks whether the work concerns \emph{opinion} (O)
    or \emph{behavior} (B). Bottom row is the gap our project fills.}
  \label{tab:rw}
  \begin{tabularx}{\textwidth}{@{}>{\raggedright\arraybackslash}p{3.0cm} >{\raggedright\arraybackslash}X c c c >{\raggedright\arraybackslash}p{2.5cm} c c@{}}
    \toprule
    \textbf{Lab work} & \textbf{Everyday medium} & \textbf{Reduce polariz.} & \textbf{LLM/ GenAI} & \textbf{Multi-agent} & \textbf{Influence lever} & \textbf{Target} & \textbf{Play} \\
    \midrule
    Ronaldo's a poser~\cite{zeng2025ronaldo}      & Online forum (sports debate) & \prt & \yes & \no  & Argument construction        & O   & \no  \\
    Climate-change TikTok~\cite{zhang2026climate} & Social media (short video)   & \no  & \no  & \no  & Personal narrative / emotion & O   & \no  \\
    ADHD on Instagram~\cite{zhang2026adhd}        & Social media (memes/comments)& \no  & \no  & \no  & Community engagement / humor & O   & \prt \\
    Online dating profiles~\cite{zhang2025dating} & Online dating platform       & \no  & \no  & \no  & Self-presentation            & O   & \no  \\
    Can AI Prompt Humans~\cite{zhang2025ecoecho}  & Serious game (NPC dialogue)  & \no  & \yes & \yes & Prompt act + show consequence& B   & \yes \\
    Breaking the News~\cite{tang2025breaking}     & Social-media-sim game        & \prt & \yes & \yes & Role-play + inoculation      & O   & \yes \\
    Eternagram~\cite{zhou2024eternagram}          & Chat / text adventure        & \prt & \yes & \no  & Narrative immersion          & O   & \yes \\
    Designing for the Invisible~\cite{du2025invisible} & Serious game            & \no  & \yes & \no  & Empathy / perspective-taking & O   & \yes \\
    Hear You in Silence~\cite{jiang2026hear}      & 1:1 conversational agent     & \no  & \yes & \no  & Active listening / pacing    & \prt& \no  \\
    \midrule
    \textbf{Ours (Polariz. Defusing)} & \textbf{Group chat / forum / Discord} & \textbf{\yes} & \textbf{\yes} & \textbf{\yes} & \textbf{Peer social influence} & \textbf{O\,/\,B} & \textbf{\prt} \\
    \bottomrule
  \end{tabularx}
\end{table*}

% =====================================================================
\subsection{Opinion Expression and Community Dynamics on Everyday Platforms}
\label{sec:rw-polarization}

Our group has studied how people voice opinions and coalesce into groups across
the very media our project targets. Most directly, \emph{Ronaldo's a poser!}~\cite{zeng2025ronaldo}
observes soccer fans debating in a live online forum, showing how participants
recruit generative AI to construct arguments and how their rhetoric shifts as
new, opposing members arrive---a close model of the competitive, identity-driven
dynamics that harden into polarization. Complementing this, work on how creators
express emotions, opinions, and data about climate change in TikTok short
videos~\cite{zhang2026climate} and on how an ADHD community engages through memes
and comments on Instagram~\cite{zhang2026adhd} characterizes opinion and
affective expression in everyday social media, including the humor and personal
narrative through which contested or stigmatized stances circulate. At the level
of identity, our study of how online-dating users construct profiles for
effective self-presentation~\cite{zhang2025dating} shows how people perform and
manage identity for an imagined audience---an antecedent of the in-group
signaling that fuels polarization. Together these studies establish the
\emph{problem space}: large groups on everyday platforms express opinion in
emotionally charged, identity-laden ways, motivating our choice to recruit
already-polarized communities and to treat polarization as both an opinion and an
affective phenomenon. What none of this work attempts is to \emph{intervene} in
those dynamics, which is the project's aim.

% =====================================================================
\subsection{LLM and Multi-Agent Systems for Social Influence and Behavior Change}
\label{sec:rw-multiagent}

A second strand of our work builds AI systems that actively steer opinion or
behavior. \emph{Can AI Prompt Humans?}~\cite{zhang2025ecoecho} inverts the usual
human-prompts-AI relationship: multimodal agents prompt players' actions and
surface their consequences, significantly shifting intended behavior---direct
evidence that a \emph{population} of agents can initiate influence rather than
merely respond, and the closest precedent for our multi-agent design.
\emph{Breaking the News}~\cite{tang2025breaking} uses LLM agents to simulate a
reactive public whose ``opinion'' players try to sway, demonstrating that
LLM-driven social audiences can be believable and consequential---the machinery
a defusing intervention would need to model peer reaction at scale. At a more
foundational level, our analysis of moral responses in large language
models~\cite{hui2025moral} quantifies how LLMs encode attitudes and value
judgments, informing both how agents might be steered to model moderate
positions and where their biases could distort a depolarization attempt. These
projects show our group can make GenAI agents that influence people and that
simulate social others, but they pursue sustainability or misinformation goals,
not the reduction of interpersonal polarization, and only \emph{Can AI Prompt
Humans} uses more than one agent.

% =====================================================================
\subsection{Conversational Agents and Playful Interventions for Attitude Change}
\label{sec:rw-intervention}

The third strand concerns how an agent should engage so that polarized users
stay in dialogue rather than entrench. \emph{Hear You in Silence}~\cite{jiang2026hear}
designs conversational agents that perform active listening through context-aware
pacing and strategic silence, showing that \emph{how} and \emph{when} an agent
responds shapes human-likeness, trust, and self-disclosure---the very levers a
defusing agent needs to earn standing before it ever challenges a view. On the
playful side, our group has repeatedly used games and speculative play as a
low-reactance route to attitude change: \emph{Eternagram}~\cite{zhou2024eternagram}
probes and nudges climate attitudes through a ChatGPT-driven text adventure;
\emph{Designing for the Invisible}~\cite{du2025invisible} uses an LLM-based
serious game to build empathy and perspective-taking toward a stigmatized
out-group; and \emph{Cracking Aegis}~\cite{fu2025aegis} turns adversarial
role-play with an LLM into awareness change. \emph{Bricolage}~\cite{agcal2025bricolage}
further shows that playful, participatory speculation can align people with a
contested cause. These designs establish that our group can shift attitudes
through conversation and play---grounding the optional \emph{play} dimension of
our project---yet each addresses a single user or a non-adversarial topic, rather
than defusing entrenched disagreement among a polarized group.

% =====================================================================
\paragraph{Summary and gap.}
Across these three strands, our group has separately established that people
polarize through emotionally charged, identity-driven expression on everyday
platforms (\S\ref{sec:rw-polarization}); that our LLM and multi-agent systems can
steer opinion and behavior and credibly simulate social others
(\S\ref{sec:rw-multiagent}); and that our conversational and playful designs can
shift attitudes while keeping users engaged (\S\ref{sec:rw-intervention}). As
Table~\ref{tab:rw} shows, no single prior project combines these commitments:
intervening in an \emph{everyday} group medium, targeting \emph{already-polarized}
populations, and using a \emph{multi-agent} community of GenAI peers to defuse
\emph{either} opinion or behavioral polarization through social influence. Our
project unites these threads to occupy precisely that gap.
.
% Citations use natbib (\citet / \cite); keys live in refs.bib.
% =====================================================================

\section{Introduction}
\label{sec:intro}

Polarization has become a defining feature of everyday online life. On the
platforms where people now spend ordinary social time---group chats, online
forums, and Discord-style channels---disagreement frequently hardens into
hostility rather than dialogue. This is not merely a matter of people holding
different views: polarization online is often \emph{affective} and
identity-driven, such that interactions between ideologically opposed users are
systematically more negative than those among like-minded
ones~\citep{marchal2022polarization}, and contact between rival groups can
amplify rather than soften antagonism~\citep{efstratiou2022nonpolar,
zhang2019sportshate}. Within tightly bonded communities, emotionally charged
expression escalates into aggression, trolling, and retaliatory spirals that
entrench positions and drive moderate voices out~\citep{wang2023emotional}. The
result is a familiar pathology of large interpersonal groups: the more people
talk, the further apart they move.

Existing responses are poorly matched to this dynamic. Content moderation,
fact-checking, and top-down corrective messaging treat polarization as a
problem of bad information delivered to individuals, yet people actively
recognize and resist perceived persuasion attempts, so overt correction tends
to provoke reactance in exactly the entrenched users it
targets~\citep{friestad1994persuasion}. These interventions also operate one
user at a time, against a phenomenon that is fundamentally \emph{social}---
sustained by group norms, peer signaling, and the felt presence of an audience.
A growing body of work shows that generative AI is now deeply entangled in these
group dynamics, with users recruiting large language models (LLMs) to sharpen
arguments and win debates in live forums~\citep{zeng2025ronaldo}. If generative
AI can amplify polarized debate, the inverse question follows naturally: can a
\emph{population} of AI agents instead be designed to defuse it?

We propose \emph{polarization defusing}: intervening directly in the everyday
chat medium of an already-polarized community with a population of
generative-AI agents that exert lateral \emph{social influence} rather than
top-down moderation. Two recent capabilities make this feasible. First,
generative agents can credibly simulate human behavior at scale---maintaining
memory, reflection, and plausible social action---making a community of LLM
agents a viable substrate for staging influence \emph{inside} a group rather
than addressing users individually~\citep{park2023generative}. Second, AI
interlocutors positioned as peers and mediators, rather than authorities, can
broaden perspectives and depolarize audiences in online
debate~\citep{govers2024aimediators}. Our central design bet is that a chorus of
peer-like agents---modeling moderate positions, normalizing disagreement, and
listening before they challenge---can reduce the salience of the persuasion
attempt itself and thereby move people where a single authoritative debunker
cannot.

Concretely, we envision a forum or Discord-style channel seeded with multiple
GenAI agents and populated with recruited members of a polarized group. The
agents draw on conversational strategies known to build the trust that makes
attitude change possible---active listening, backchanneling, and reciprocal
self-disclosure~\citep{xiao2020ifihear, jiang2026hear}---and may optionally
employ light playful or inoculation-style framings that make encountering an
opposing view feel exploratory rather than confrontational~\citep{roozenbeek2019badnews}.
This design lets us study defusing along two axes that prior work has treated
separately: whether the intervention shifts \emph{opinion} (expressed stance,
affect toward the out-group) or \emph{behavior} (how members post, escalate, or
engage). Accordingly, we ask:
\begin{itemize}
  \item \textbf{RQ1.} Can a multi-agent population of GenAI peers reduce
  affective and opinion polarization within an already-polarized online group,
  relative to no intervention or a single-agent mediator?
  \item \textbf{RQ2.} Which mechanisms of social influence---peer modeling,
  active listening, normalizing disagreement, playful reframing---drive
  defusing, and how do they interact?
  \item \textbf{RQ3.} How do effects on expressed \emph{opinion} relate to
  effects on group \emph{behavior}, and do they persist beyond the
  intervention?
\end{itemize}

This work makes the following anticipated contributions: (1)~the framing and
design of \emph{polarization defusing} as a multi-agent social-influence
intervention situated in an everyday group medium; (2)~a system that embeds
peer-like GenAI agents with active-listening and playful strategies in a
forum/Discord-style channel; and (3)~empirical findings on whether and how such
agents defuse opinion and behavioral polarization in recruited, already-polarized
groups. The remainder of the paper situates these contributions across three
lines of prior work---how online groups polarize, how LLM and multi-agent
systems carry social influence, and how conversational and playful interventions
shift attitudes---which we review next.
